% Part: incompleteness
% Chapter: introduction
% Section: definitions

\documentclass[../../../include/open-logic-section]{subfiles}

\begin{document}

\olfileid{inc}{int}{def}

\olsection{সংজ্ঞা}

গণিতকে বিধিবদ্ধ করা এবং এমন বিধিবদ্ধ রূপ সঙ্গত ও পূর্ণ—এটি দেখানোর হিলবার্টীয়
কর্মসূচি বাস্তবায়নের প্রথম কাজ হলো একটি ভাষা, যৌক্তিক কাঠামো ও স্বতঃসিদ্ধের
ব্যবস্থা বেছে নেওয়া। আমাদের উদ্দেশ্যে ধরে নিই, গণিতকে প্রথম-ক্রমের একটি ভাষায়
বিধিবদ্ধ করা যায়; অর্থাৎ !!{constant}s, !!{function}s ও !!{predicate}s-এর এমন
একটি সেট আছে যা প্রথম-ক্রমের যুক্তিবিদ্যার সংযোজক ও পরিমাণসূচকের সঙ্গে মিলে
গণিতের দাবিগুলি প্রকাশ করতে দেয়। অধিকাংশ মানুষই একমত যে এমন ভাষা আছে: সেটতত্ত্বের
ভাষা, যেখানে একমাত্র অযৌক্তিক প্রতীক~$\in$। এত সরল ভাষা এত বেশি প্রকাশক্ষম—প্রথম
দর্শনে দাবিটি অবশ্যই খুব অবিশ্বাস্য; এই অতি সংযত শব্দভাণ্ডারে কার্যত সমগ্র গণিত
প্রকাশ করা যায়, তা প্রতিষ্ঠায় অনেক কাজ লেগেছে। আপাতত বিষয়টি সরল রাখতে আলোচনা
পাটিগণিতে সীমিত করি, অর্থাৎ গণিতের যে অংশ কেবল স্বাভাবিক সংখ্যা~$\Nat$ নিয়ে
কাজ করে। পাটিগণিতের তথ্য প্রকাশের স্বাভাবিক ভাষা~$\Lang L_A$। $\Lang L_A$-তে
আছে একটি দ্বি-স্থানীয় !!{predicate}~$<$, একটি !!{constant}~$\Obj 0$, একটি
এক-স্থানীয় !!{function}~$\prime$, এবং দুটি দ্বি-স্থানীয় !!{function}s~$+$
ও~$\times$।

\begin{defn}
!!{sentence}s-এর একটি সেট~$\Gamma$ অনুগমনের অধীনে আবদ্ধ হলে তাকে একটি
\emph{তত্ত্ব} বলা হয়; অর্থাৎ $\Gamma = \Setabs{!A}{\Gamma \Entails
!A}$ হলে।
\end{defn}

তত্ত্ব নির্দিষ্ট করার দুটি সহজ উপায় আছে। একটি হলো, কোনো !!{structure}-তে সত্য
!!{sentence}s-গুলির সেট হিসেবে দেওয়া। যেমন, $\Lang L_A$-এর সেই
!!{structure}-টি বিবেচনা করো যার !!{domain} হলো~$\Nat$ এবং যাতে সব অযৌক্তিক
প্রতীকের প্রত্যাশিত ব্যাখ্যা দেওয়া হয়।

\begin{defn}\ollabel{def:standard-model}
\emph{পাটিগণিতের মানক মডেল} হলো নিচের মতো সংজ্ঞায়িত !!{structure}~$\Struct
N$:
\begin{enumerate}
\item $\Domain N = \Nat$
\item $\Assign{\Obj 0}{N} = 0$
\item সব $n \in \Nat$-এর জন্য $\Assign{\Obj \prime}{N}(n) = n + 1$
\item সব $n, m \in \Nat$-এর জন্য $\Assign{\Obj +}{N}(n, m) = n + m$
\item সব $n, m \in \Nat$-এর জন্য $\Assign{\Obj \times}{N}(n, m) = n\cdot m$
\item $\Assign{\Obj <}{N} = \Setabs{\tuple{n, m}}{n \in \Nat, m \in
  \Nat, n < m}$
\end{enumerate}
\end{defn}

`$\times$' ও `$\cdot$'-এর পার্থক্য লক্ষ করো: $\times$ পাটিগণিতের ভাষার একটি
প্রতীক। গুণনের কথা মনে করিয়ে দেওয়ার জন্যই প্রতীকটি বেছে নেওয়া হয়েছে, কিন্তু
$\times$ গুণন-ক্রিয়া নয়; এটি একটি দ্বি-স্থানীয় অপেক্ষক-প্রতীক (আনুষ্ঠানিকভাবে
$\Obj f^2_1$)। অন্যদিকে, $\cdot$-ই~\emph{হলো} সাধারণ গুণন-অপেক্ষক। $n \cdot m$
জাতীয় কিছু দেখলে তার অর্থ $n$ ও~$m$ সংখ্যার গুণফল; আর $x \times y$ জাতীয়
কিছু দেখলে পাটিগণিতের ভাষার একটি পদ নিয়ে কথা হচ্ছে। মানক মডেলে অপেক্ষক-প্রতীক
$\times$-এর ব্যাখ্যা হলো স্বাভাবিক সংখ্যার উপর অপেক্ষক~$\cdot$। যোগের ক্ষেত্রে
পাটিগণিতের ভাষার অপেক্ষক-প্রতীক এবং স্বাভাবিক সংখ্যার যোগ-অপেক্ষক—উভয়ের জন্যই
আমরা $+$ ব্যবহার করি। এখানে প্রসঙ্গ দেখে কোন অর্থটি অভিপ্রেত তা বুঝতে হবে।

\begin{defn}
\emph{সত্য পাটিগণিতের} তত্ত্ব হলো পাটিগণিতের মানক মডেলে পরিতৃপ্ত
!!{sentence}s-গুলির সেট; অর্থাৎ
\[
\Th{TA} = \Setabs{!A}{\Sat{N}{!A}}.
\]
\end{defn}

$\Th{TA}$ একটি তত্ত্ব। কারণ $\Th{TA} \Entails !A$ হলেই $\Th{TA}$-কে
পরিতৃপ্ত করে এমন প্রতিটি !!{structure}-তে $!A$ পরিতৃপ্ত। যেহেতু
$\Sat{N}{\Th{TA}}$, তাই $\Sat{N}{!A}$; অতএব $!A \in \Th{TA}$।

তত্ত্ব~$\Gamma$ নির্দিষ্ট করার অন্য উপায় হলো কোনো বাক্যসমষ্টি~$\Gamma_0$ থেকে
অনুগত !!{sentence}s-গুলির সেট হিসেবে দেওয়া। সে ক্ষেত্রে $\Gamma$ হলো অনুগমনের
অধীনে $\Gamma_0$-এর ``আবরণ''। এইভাবে তত্ত্ব নির্দিষ্ট করা তখনই আগ্রহজনক যখন
$\Gamma_0$ স্পষ্টভাবে নির্দিষ্ট—যেমন, $\Gamma_0$-এর !!{element}s-গুলি তালিকাভুক্ত
থাকে। অন্ততপক্ষে $\Gamma_0$ নির্ণেয় হওয়া চাই; অর্থাৎ কোনো !!a{sentence}
$\Gamma_0$-এর উপাদান কি না তার একটি গণনসাধ্য পরীক্ষা থাকতে হবে। $\Gamma_0$-এর
!!{sentence}s-গুলিকে আমরা $\Gamma$-এর \emph{স্বতঃসিদ্ধ} বলি, এবং বলি
$\Gamma_0$ দ্বারা $\Gamma$ \emph{স্বতঃসিদ্ধায়িত}।

\begin{defn}
তত্ত্ব~$\Gamma$ $\Gamma_0$ দ্বারা \emph{স্বতঃসিদ্ধায়িত} যদি এবং কেবল যদি
\[
\Gamma = \Setabs{!A}{\Gamma_0 \Entails !A}
\]
\end{defn}

\begin{defn}
নিচের বাক্যগুলি দ্বারা স্বতঃসিদ্ধায়িত তত্ত্ব~$\Th{Q}$ ``রবিনসনের
$\Th{Q}$'' নামে পরিচিত; এটি পাটিগণিতের একটি অতি সরল তত্ত্ব।
\begin{align*}
& \lforall[x][\lforall[y][(\eq[x'][y'] \lif \eq[x][y])]] \tag{$!Q_1$}\\
& \lforall[x][\eq/[\Obj 0][x']] \tag{$!Q_2$}\\
& \lforall[x][(\eq[x][\Obj 0] \lor \lexists[y][\eq[x][y']])] \tag{$!Q_3$}\\
& \lforall[x][\eq[(x + \Obj 0)][x]] \tag{$!Q_4$}\\
& \lforall[x][\lforall[y][\eq[(x + y')][(x + y)']]] \tag{$!Q_5$}\\
& \lforall[x][\eq[(x \times \Obj 0)][\Obj 0]] \tag{$!Q_6$}\\
& \lforall[x][\lforall[y][\eq[(x \times y')][((x \times y) + x)]]] \tag{$!Q_7$}\\
& \lforall[x][\lforall[y][(x < y \liff \lexists[z][\eq[(z' + x)][y]])]] \tag{$!Q_8$}
\end{align*}
!!{sentence}s-সমষ্টি $\{!Q_1, \dots, !Q_8\}$ হলো $\Th{Q}$-এর স্বতঃসিদ্ধ;
তাই এগুলি থেকে অনুগত সব !!{sentence}s নিয়ে $\Th{Q}$ গঠিত:
\[
\Th{Q} = \Setabs{!A}{\{!Q_1, \dots, !Q_8\} \Entails !A}.
\]
\end{defn}

\begin{defn}
ধরা যাক, $!A(x)$ হলো $\Lang L_A$-এর এমন একটি !!a{formula} যার মুক্ত চলরাশি
$x$ এবং $y_1$, \dots, $y_n$। তাহলে নিচের রূপের প্রতিটি !!{sentence} হলো
\emph{আরোহ-ছকের} একটি নিদর্শন:
\[
\lforall[y_1][\dots\lforall[y_n][((!A(\Obj 0) \land \lforall[x][(!A(x)
\lif !A(x'))]) \lif \lforall[x][!A(x)])]]
\]

\emph{পেয়ানো পাটিগণিত}~$\Th{PA}$ হলো $\Th{Q}$-এর স্বতঃসিদ্ধ এবং আরোহ-ছকের
সব নিদর্শন দিয়ে স্বতঃসিদ্ধায়িত তত্ত্ব।
\end{defn}

\begin{explain}
আরোহ-ছকের প্রতিটি নিদর্শন~$\Struct{N}$-এ সত্য। !!{formula}~$!A$-এর একমাত্র
মুক্ত !!{variable}~$x$ হলে এটি সবচেয়ে সহজে বোঝা যায়। তখন $!A(x)$
$\Struct{N}$-এ $\Nat$-এর একটি উপসেট~$X_{!A}$ সংজ্ঞায়িত করে। $X_{!A}$ হলো
সেই সব $n \in \Nat$-এর সেট, যাদের জন্য $s(x) = n$ হলে
$\Sat{N}{!A(x)}[s]$। আরোহ-ছকের সংশ্লিষ্ট নিদর্শনটি হলো
\[
((!A(\Obj 0) \land \lforall[x][(!A(x) \lif !A(x'))]) \lif
  \lforall[x][!A(x)]).
\]
এর পূর্বপক্ষ~$\Struct{N}$-এ সত্য হলে $0 \in X_{!A}$; আর যখনই
$n \in X_{!A}$, তখন $n+1$-ও এতে থাকে। $0 \in X_{!A}$ বলে পাই
$1 \in X_{!A}$। $1 \in X_{!A}$ থেকে পাই $2 \in X_{!A}$। এভাবে চলতে
থাকে। তাই প্রতিটি $n \in \Nat$-এর জন্য $n \in X_{!A}$। কিন্তু এর অর্থ
$\lforall[x][!A(x)]$ বাক্যটি~$\Struct{N}$-এ পরিতৃপ্ত।
\end{explain}

$\Th{Q}$ ও $\Th{PA}$ উভয়ই স্বতঃসিদ্ধায়িত তত্ত্ব। বড় প্রশ্ন হলো, এরা কতটা
শক্তিশালী? যেমন, $\Lang L_A$-এ প্রকাশযোগ্য~$\Nat$-সম্বন্ধীয় সব সত্য কি
$\Th{PA}$ প্রমাণ করতে পারে? নির্দিষ্টভাবে, $\Lang L_A$-এ সূত্রায়িত করা যায়
এমন সব প্রশ্ন কি $\Th{PA}$-এর স্বতঃসিদ্ধগুলি মীমাংসা করে?

প্রশ্নটি অন্যভাবেও করা যায়: $\Th{PA} = \Th{TA}$ কি? $\Th{TA}$ স্পষ্টতই
$\Nat$-সম্বন্ধীয় সব সত্য প্রমাণ করে (অর্থাৎ তাদের অন্তর্ভুক্ত করে), এবং
$\Lang L_A$-এ সূত্রায়িত করা যায় এমন সব প্রশ্ন মীমাংসা করে। কারণ $!A$
$\Lang L_A$-এর কোনো !!a{sentence} হলে হয় $\Sat{N}{!A}$, নয়
$\Sat{N}{\lnot !A}$; অতএব হয় $\Th{TA} \Entails !A$, নয়
$\Th{TA} \Entails \lnot !A$। এমন তত্ত্বকে \emph{!!{complete}} বলি।

\begin{defn}
তত্ত্ব $\Gamma$ \emph{!!{complete}} যদি এবং কেবল যদি তার ভাষার প্রতিটি
!!{sentence}~$!A$-এর জন্য হয় $\Gamma \Entails !A$, নয়
$\Gamma \Entails \lnot !A$।
\end{defn}

\begin{explain}
পূর্ণতা উপপাদ্য অনুসারে $\Gamma \Entails !A$ যদি এবং কেবল যদি $\Gamma \Proves
!A$। তাই $\Gamma$ পূর্ণ যদি এবং কেবল যদি তার ভাষার প্রতিটি !!{sentence}~$!A$-এর
জন্য হয় $\Gamma \Proves !A$, নয় $\Gamma \Proves \lnot !A$।
\end{explain}

এ থেকে আরেকটি প্রশ্ন আসে: কোনো !!a{sentence}~$\Th{TA}$, $\Th{PA}$, এমনকি শুধু
$\Th{Q}$-এর অন্তর্ভুক্ত কি না পরীক্ষা করার কোনো গণনপদ্ধতি কি আছে? কোনো সেট
(যেমন !!{sentence}s-এর সেট) কখন !!{decidable}, তা সংজ্ঞায়িত করে প্রশ্নটিকে আরও
নির্দিষ্ট করা যায়।

\begin{defn}
একটি সেট $X$ \emph{!!{decidable}} যদি এবং কেবল যদি এমন একটি গণনপদ্ধতি থাকে
যা নিবেশ~$x$-এ $x \in X$ হলে~$1$, অন্যথায়~$0$ ফেরত দেয়।
\end{defn}

তাহলে আমাদের প্রশ্ন দাঁড়ায়: $\Th{TA}$ ($\Th{PA}$, $\Th{Q}$) কি !!{decidable}?

এই সব প্রশ্নের উত্তর হবে: না। এই তত্ত্বগুলির কোনোটিই নির্ণেয় নয়। তবে ঘটনাটি
শুধু এই বিশেষ তত্ত্বগুলির বৈশিষ্ট্য নয়। বস্তুত, নির্দিষ্ট কয়েকটি শর্ত পূরণকারী
\emph{প্রতিটি} তত্ত্বের ক্ষেত্রেই একই ফল প্রযোজ্য। $\Th{Q}$ ও $\Th{PA}$ যে
শর্তগুলির একটি পূরণ করে তা হলো: তত্ত্বগুলি !!a{decidable} স্বতঃসিদ্ধ-সমষ্টি
দ্বারা স্বতঃসিদ্ধায়িত।

\begin{defn}
কোনো তত্ত্ব !!a{decidable} স্বতঃসিদ্ধ-সমষ্টি দ্বারা স্বতঃসিদ্ধায়িত হলে তাকে
\emph{!!{axiomatizable}} বলা হয়।
\end{defn}

\begin{ex}
!!{sentence}s-এর কোনো সসীম সেট দ্বারা স্বতঃসিদ্ধায়িত প্রতিটি তত্ত্বই
!!{axiomatizable}, কারণ প্রতিটি সসীম সেট !!{decidable}। তাই, যেমন,
$\Th{Q}$ !!{axiomatizable}।

$\Th{PA}$-এর মতো ছক দ্বারা স্বতঃসিদ্ধায়িত তত্ত্বও !!{axiomatizable}। কোনো
$!B$ $\Th{PA}$-এর স্বতঃসিদ্ধ কি না পরীক্ষা করতে—অর্থাৎ $!B$ $\Th{PA}$-এর
স্বতঃসিদ্ধ হলে $\Char{X}(!B) = 1$, অন্যথায় $= 0$ এমন অপেক্ষক~$\Char{X}$
গণনা করতে—আমরা নিচের কাজ করতে পারি। প্রথমে পরীক্ষা করো $!B$ কি $\Th{Q}$-এর
কোনো স্বতঃসিদ্ধ। হলে উত্তর ``হ্যাঁ'' এবং $\Char{X}(!B) = 1$। না হলে পরীক্ষা
করো এটি আরোহ-ছকের কোনো নিদর্শন কি না। এটি সুসংবদ্ধভাবে করা যায়; এখানে সম্ভবত
নিচেরভাবে তা সবচেয়ে সহজে বোঝা যায়। আরোহ-ছকের প্রতিটি নিদর্শনের শুরুতে কয়েকটি
সার্বিক পরিমাণসূচক থাকে, তারপর থাকে একটি প্রকল্পবিশিষ্ট উপ-!!{formula}। সেই
প্রকল্পটির অনুবর্তী হলো $\lforall[x][!A(x, y_1, \dots, y_n)]$, যেখানে $x$ ও
$y_1$, \dots, $y_n$ হলো $!A$-এর সব মুক্ত চলরাশি, এবং $!B$-এর গোড়ার
পরিমাণসূচকগুলি $y_1$, \dots,~$y_n$ চলরাশিগুলিকে বাঁধে। এই~$!A$ নিষ্কাশন করে
তার মুক্ত চলরাশিগুলি গোড়ার সার্বিক পরিমাণসূচক ও~$\lforall[x]$ দ্বারা বাঁধা
চলরাশিগুলির সঙ্গে মেলে কি না পরীক্ষা করার পর আমরা দেখি প্রকল্পটির পূর্বপক্ষ
নিচেরটির সঙ্গে মেলে কি না:
\[
!A(\Obj 0, y_1, \dots, y_n) \land \lforall[x][(!A(x, y_1, \dots, y_n)
\lif !A(x', y_1, \dots, y_n))]
\]
আবার, মিললে $!B$ আরোহ-ছকের একটি নিদর্শন, না মিললে $!B$ তা নয়।
% BN-SRC-204 TARGET-CORRECTION-BEGIN
\par\smallskip\noindent\textnormal{\small\textbf{উৎস-সংশোধন (BN-SRC-204):} হিমায়িত অনুচ্ছেদের শেষ বাক্যে দ্বিতীয় শাখায় $!B$-এর উল্লেখটি হারিয়ে গিয়েছিল। বাংলা পাঠে ``না মিললে $!B$ তা নয়'' বলে একই বাক্যের উভয় শাখায় পরীক্ষিত সূত্রটি স্পষ্ট রাখা হয়েছে।} % BN-SRC-204 TARGET-CORRECTION-END
\end{ex}

এই প্রশ্নের উত্তর দিতে—এবং কোন তত্ত্ব পূর্ণ বা নির্ণেয় সেই আরও সাধারণ প্রশ্নের
উত্তর দিতে—নিচের সংজ্ঞাটিও কাজে লাগবে। মনে করো, সেট $X$ !!{enumerable} যদি
এবং কেবল যদি সেটি শূন্য হয় অথবা কোনো !!a{surjective} অপেক্ষক~$f \colon \Nat
\to X$ থাকে। এমন অপেক্ষককে $X$-এর একটি তালিকায়ন বলা হয়।

\begin{defn}
একটি সেট $X$-কে \emph{!!{computably enumerable}} (সংক্ষেপে !!{c.e.}) বলা হয়
যদি এবং কেবল যদি সেটি শূন্য অথবা তার একটি গণনসাধ্য তালিকায়ন থাকে।
\end{defn}

!!{axiomatizability}-র পাশাপাশি অসম্পূর্ণতা উপপাদ্য প্রযোজ্য হওয়ার আরেকটি শর্ত
হলো, তত্ত্বগুলি গণনসাধ্য অপেক্ষক ও !!{decidable} সম্বন্ধের মৌলিক তথ্য প্রমাণের
পক্ষে যথেষ্ট শক্তিশালী। ``মৌলিক তথ্য'' বলতে আমরা এমন !!{sentence}s বুঝি যা
গণনসাধ্য অপেক্ষকের প্রতিটি আর্গুমেন্টে তার মান প্রকাশ করে। আর ``যথেষ্ট
শক্তিশালী'' বলতে বুঝি, সংশ্লিষ্ট তত্ত্ব এই বাক্যগুলিকে তার উপপাদ্য হিসেবে গণ্য
করে। যেমন, একটি আদর্শ গণনসাধ্য অপেক্ষক যোগ বিবেচনা করো। $2$ ও $3$ আর্গুমেন্টে
$+$-এর মান~$5$; অর্থাৎ $2+3 = 5$। $2$ ও $3$ আর্গুমেন্টে $+$-এর মান~$5$—এই
তথ্য প্রকাশকারী পাটিগণিতের ভাষার একটি বাক্য হলো
$(\num{2} + \num{3}) = \num{5}$। যেমন, $\Th{Q}$ এই বাক্যটি প্রমাণ করে। আরও
সাধারণভাবে, আমরা চাই প্রতিটি গণনসাধ্য অপেক্ষক $f(x_1, x_2)$-এর জন্য
$\Lang{L_A}$-এ এমন একটি !!a{formula}~$!A_f(x_1, x_2, y)$ থাকুক, যাতে যখনই
$f(n_1, n_2) = m$, তখন $\Th{Q} \Proves !A_f(\num{n_1}, \num{n_2},
\num{m})$। এইভাবে $\Th{Q}$ প্রমাণ করে যে $n_1$, $n_2$ আর্গুমেন্টে~$f$-এর
মান~$m$। বস্তুত আমরা চাই এটি আরও একটু বেশি প্রমাণ করুক: $n_1$,~$n_2$
আর্গুমেন্টে অন্য কোনো সংখ্যা~$f$-এর মান নয়। !!{decidable} সম্বন্ধের ক্ষেত্রেও
একই কথা। নিচের দুই সংজ্ঞায় তা নির্দিষ্ট করা হলো।

\begin{defn}
একটি !!{formula}~$!A(x_1, \dots, x_k, y)$ তত্ত্ব~$\Gamma$-তে অপেক্ষক
$f\colon \Nat^k \to \Nat$-কে \emph{!!{represents}} করে যদি এবং কেবল যদি
যখনই $f(n_1, \dots, n_k) = m$, তখন
\begin{enumerate}
\item $\Gamma \Proves !A(\num{n_1}, \dots, \num{n_k}, \num{m})$, এবং
\item $\Gamma \Proves \lforall[y](!A(\num{n_1}, \dots, \num{n_k},
y) \lif y = \num{m})$।
\end{enumerate}
\end{defn}

\begin{defn}
একটি !!{formula}~$!A(x_1, \dots, x_k)$ প্রদত্ত তত্ত্বে সম্বন্ধ
$R \subseteq \Nat^k$-কে \emph{!!{represents}} করে যদি এবং কেবল যদি,
\begin{enumerate}
\item যখনই $R(n_1, \dots, n_k)$, তখন $\Gamma \Proves !A(\num{n_1},
\dots, \num{n_k})$, এবং
\item যখনই $R(n_1, \dots, n_k)$ নয়, তখন $\Gamma \Proves \lnot
!A(\num{n_1}, \dots, \num{n_k})$।
\end{enumerate}
% BN-SRC-200 TARGET-CORRECTION-BEGIN
\par\smallskip\noindent\textnormal{\small\textbf{উৎস-সংশোধন (BN-SRC-200):} হিমায়িত সংজ্ঞার শিরোনামে অপেক্ষকের সংজ্ঞাটির মতো তত্ত্বের পরামিতি উল্লেখ করা হয়নি, অথচ উভয় শর্তেই $\Gamma$-তে প্রমাণযোগ্যতা ব্যবহৃত। বাংলা পাঠে ``প্রদত্ত তত্ত্বে'' যোগ করে এই নির্ভরতা স্পষ্ট করা হয়েছে।} % BN-SRC-200 TARGET-CORRECTION-END
\end{defn}

অসম্পূর্ণতা উপপাদ্য প্রয়োগের পক্ষে একটি তত্ত্ব তখনই ``যথেষ্ট শক্তিশালী'' যখন
তা সব গণনসাধ্য অপেক্ষক ও সব !!{decidable} সম্বন্ধকে !!{represents} করে।
$\Th{Q}$ ও তার সম্প্রসারণগুলি এই শর্ত পূরণ করে, কিন্তু তা প্রতিষ্ঠা করতে আমাদের
কিছুটা সময় লাগবে। $\Th{Q}$ কী ধরনের তথ্য প্রমাণ করতে পারে, এটি সে বিষয়ে একটি
অ-তুচ্ছ তথ্য; আর $\Th{Q}$-এর অল্প কয়েকটি স্বতঃসিদ্ধ থেকেই এই সব তথ্য প্রমাণ
করতে হবে বলে তা দেখানো কঠিন। অথচ $\Th{Q}$ খুব দুর্বল তত্ত্ব। তাই $\Th{Q}$ সব
গণনসাধ্য অপেক্ষক উপস্থাপন করে—এটি প্রমাণ করা কঠিন হলেও অধিকাংশ আগ্রহজনক তত্ত্ব
$\Th{Q}$-এর চেয়ে শক্তিশালী, অর্থাৎ $\Th{Q}$-এর চেয়ে বেশি কিছু প্রমাণ করে।
আর $\Th{Q}$ কিছু প্রমাণ করলে প্রতিটি শক্তিশালী তত্ত্বও তা প্রমাণ করে; $\Th{Q}$
যেহেতু সব গণনসাধ্য অপেক্ষক উপস্থাপন করে, তাই প্রতিটি শক্তিশালী তত্ত্বও করে।
এর অর্থ বহু আগ্রহজনক তত্ত্ব অসম্পূর্ণতা উপপাদ্যের এই শর্তটি পূরণ করে। অতএব
আমাদের কঠিন পরিশ্রম সার্থক হবে, কারণ তা দেখাবে যে অসম্পূর্ণতা উপপাদ্য বহু ধরনের
তত্ত্বে প্রযোজ্য। নিশ্চয়ই ``সমগ্র গণিত'' বিধিবদ্ধ করার লক্ষ্যবিশিষ্ট কোনো
তত্ত্বকে $\Th{Q}$ যা প্রমাণ করে তার সবই প্রমাণ করতে হবে, কারণ অন্তত প্রাথমিক
গণনার ফলগুলি ধারণ করার সামর্থ্য তার থাকা চাই। তাই ``সমগ্র গণিতের'' তত্ত্ব হওয়ার
প্রতিটি প্রার্থীই এমন একটি তত্ত্ব হবে যার ক্ষেত্রে অসম্পূর্ণতা উপপাদ্য প্রযোজ্য।

\end{document}
